\documentclass[12pt]{article}
\usepackage{amsmath,amssymb}

\begin{document}
\section*{{2023 AIME I Problems/Problem 14}}
Source: AoPS Wiki

\subsection*{Solution}
This problem is, in essence, the following: A $12\times12$ coordinate grid is placed on a (flat) torus; how many loops are there that pass through each point while only moving up or right? In other words, Felix the frog starts his journey at $(0,0)$ in the coordinate plane. Every second, he jumps either to the right or up, until he reaches an $x$ - or $y$ -coordinate of $12$ . At this point, if he tries to jump to a coordinate outside the square from $(0,0)$ to $(11,11)$ , he "wraps around" and ends up at an $x$ - or $y$ - coordinate of $0$ . How many ways are there for Felix to jump on every grid point in this square, so that he ends at $(0,0)$ ? This is consistent with the construction of the flat torus as $\mathbb Z^2/12\mathbb Z^2$ (2-dimensional modular arithmetic. $(\mathbb{Z}_{12})^2$ ) Moving on, define a $\textit{path}$ from point $A$ to point $B$ to be a sequence of "up"s and "right"s that takes Felix from $A$ to $B$ . The $\textit{distance}$ from $A$ to $B$ is the length of the shortest path from $A$ to $B$ . At the crux of this problem is the following consideration: The points $A_i=(i,11-i), i\in{0,...,11}$ are pairwise equidistant, each pair having distance of $12$ in both directions. [asy] size(7cm); for (int x=0; x<12; ++x){ for (int y=0; y<12; ++y){ fill(circle((x,y),0.05));}} for (int i=0; i<12; ++i){ fill(circle((i,11-i),0.1),red);} pen p=green+dashed; path u=(3,8)--(4,8)--(4,9)--(4,10)--(4,11)--(5,11)--(5,11.5); path v=(5,-0.5)--(5,0)--(5,1)--(6,1)--(6,2)--(6,3)--(6,4)--(7,4); draw(u,p); draw(v,p); pen p=blue+dashed; path u=(4,7)--(5,7)--(5,8)--(5,9)--(5,10)--(6,10)--(6,11)--(6,11.5); path v=(6,-0.5)--(6,0)--(7,0)--(7,1)--(7,2)--(7,3)--(8,3); draw(u,p); draw(v,p); [/asy] A valid complete path then joins two $A_i$ 's, say $A_i$ and $A_j$ . In fact, a link between some $A_i$ and $A_j$ fully determines the rest of the cycle, as the path from $A_{i+1}$ must "hug" the path from $A_i$ , to ensure that there are no gaps. We therefore see that if $A_0$ leads to $A_k$ , then $A_i$ leads to $A_{i+k}$ . Only the values of $k$ relatively prime to $12$ result in solutions, though, because otherwise $A_0$ would only lead to $\{A_i:\exists n\in \mathbb Z:i\equiv kn\quad\text{mod 12}\}$ . The number of paths from $A_0$ to $A_k$ is ${12\choose k}$ , and so the answer is \[{12\choose1}+{12\choose5}+{12\choose7}+{12\choose11}=1\boxed{608}.\] Notes: - One can prove that the path from $A_{i+1}$ must "hug" the path from $A_i$ by using techniques similar to those in Solution 2. - One can count the paths as follows: To get from $A_0$ to $A_i$ , Felix takes $k$ rights and $12-k$ ups, which can be done in ${12\choose k}$ ways.

\end{document}
